On applique la formule~:
\[
    E_c=\frac{1}{2} M v^2
\]
en utilisant les unités du Système International.

La masse $M$ doit être en kilogrammes~:
\[
    M=\qty{22}{tonnes}=\qty{22e3}{\kg}=\qty{2.2e4}{\kg}.
\]

La vitesse du véhicule est aussi celle de son centre d'inertie~; elle doit être
en $\unit{\m\per\s}$~:

D'où~:
\[
    v=\qty{72}{\km\per\h}=\frac{72000}{3600}=\qty{20}{\m\per\s}.
\]

\[
    E_c=\frac{1}{2} M v^2=
    \frac{1}{2}\times\bigl(2,2 \times 10^4\bigr) \times 20^2
    \quad \Rightarrow\
    E_c=\qty{4.4e6}{\J}=\qty{4.4}{\MJ}.
\]

